\documentclass[11pt]{article}

\usepackage[margin=0.9in]{geometry}
\usepackage{amsmath,amssymb}
\usepackage{booktabs}
\usepackage{enumitem}
\usepackage{listings}
\usepackage{xcolor}
\usepackage{hyperref}
\usepackage{tikz}

\definecolor{backcolour}{rgb}{0.95,0.95,0.92}
\definecolor{codeblue}{rgb}{0.05,0.15,0.55}
\definecolor{codegreen}{rgb}{0.0,0.45,0.0}

\lstset{
    backgroundcolor=\color{backcolour},
    basicstyle=\ttfamily\small,
    keywordstyle=\color{codeblue}\bfseries,
    commentstyle=\color{codegreen},
    stringstyle=\color{purple},
    breaklines=true,
    frame=single,
    showstringspaces=false,
    tabsize=4
}

\title{CPSC 250 \\ Plant Information Example}
\author{Edward Brash}
\date{}

\begin{document}

\maketitle

\section{Program Goal}

In this example, we use inheritance to build a small plant information program.

We are given:

\begin{itemize}
    \item a base \texttt{Plant} class,
    \item a derived \texttt{Flower} class.
\end{itemize}

We want to create a list called:

\begin{lstlisting}[language=Python]
my_garden
\end{lstlisting}

that stores both \texttt{Plant} and \texttt{Flower} objects.

\section{The Base Plant Class}

A plant has:

\begin{itemize}
    \item a name,
    \item a cost.
\end{itemize}

\begin{lstlisting}[language=Python]
class Plant:

    def __init__(self, plant_name, plant_cost):
        self.plant_name = plant_name
        self.plant_cost = plant_cost

    def print_info(self):
        print(f"Plant name: {self.plant_name}")
        print(f"Plant cost: {self.plant_cost}")
\end{lstlisting}

\section{The Derived Flower Class}

A flower is a plant, but it has additional information:

\begin{itemize}
    \item whether it is annual,
    \item flower color.
\end{itemize}

\begin{lstlisting}[language=Python]
class Flower(Plant):

    def __init__(self, plant_name, plant_cost,
                 is_annual, color_of_flowers):

        Plant.__init__(self, plant_name, plant_cost)

        self.is_annual = is_annual
        self.color_of_flowers = color_of_flowers

    def print_info(self):
        Plant.print_info(self)
        print(f"Annual: {self.is_annual}")
        print(f"Color of flowers: {self.color_of_flowers}")
\end{lstlisting}

\section{The Input Format}

The program repeatedly reads user input until the user enters:

\begin{lstlisting}
-1
\end{lstlisting}

Each input line tells us whether the item is a plant or a flower.

Examples:

\begin{lstlisting}
plant fern 12.99
flower rose 15.99 true red
-1
\end{lstlisting}

\section{Processing the Input}

We split the user input into tokens:

\begin{lstlisting}[language=Python]
tokens = user_string.split()
\end{lstlisting}

Then:

\begin{lstlisting}[language=Python]
type = tokens[0]
name = tokens[1]
cost = tokens[2]
\end{lstlisting}

If the type is \texttt{plant}, we create a \texttt{Plant} object.

If the type is \texttt{flower}, we create a \texttt{Flower} object.

\section{Creating Objects}

\begin{lstlisting}[language=Python]
if type == "plant":

    my_plant = Plant(name, cost)
    my_garden.append(my_plant)

elif type == "flower":

    is_annual = tokens[3]
    color = tokens[4]

    my_flower = Flower(name, cost, is_annual, color)
    my_garden.append(my_flower)
\end{lstlisting}

\section{The Main Loop}

\begin{lstlisting}[language=Python]
my_garden = []

user_string = input()

while user_string != "-1":

    tokens = user_string.split()

    type = tokens[0]
    name = tokens[1]
    cost = tokens[2]

    if type == "plant":

        my_plant = Plant(name, cost)
        my_garden.append(my_plant)

    elif type == "flower":

        is_annual = tokens[3]
        color = tokens[4]

        my_flower = Flower(name, cost, is_annual, color)
        my_garden.append(my_flower)

    user_string = input()
\end{lstlisting}

\section{Printing the List}

We write a function:

\begin{lstlisting}[language=Python]
def print_list(my_garden):
\end{lstlisting}

that loops over every object in the list and calls its \texttt{print\_info()} method.

\begin{lstlisting}[language=Python]
def print_list(my_garden):

    for i in range(len(my_garden)):

        print(f"Plant {i + 1} Information:")
        my_garden[i].print_info()
        print()
\end{lstlisting}

\section{Polymorphism}

This example demonstrates polymorphism.

The list contains both:

\begin{itemize}
    \item \texttt{Plant} objects,
    \item \texttt{Flower} objects.
\end{itemize}

But the loop simply calls:

\begin{lstlisting}[language=Python]
my_garden[i].print_info()
\end{lstlisting}

Python automatically uses the correct version of \texttt{print\_info()} depending on the actual object type.

\section{Complete Program}

\begin{lstlisting}[language=Python]
class Plant:

    def __init__(self, plant_name, plant_cost):
        self.plant_name = plant_name
        self.plant_cost = plant_cost

    def print_info(self):
        print(f"Plant name: {self.plant_name}")
        print(f"Plant cost: {self.plant_cost}")


class Flower(Plant):

    def __init__(self, plant_name, plant_cost,
                 is_annual, color_of_flowers):

        Plant.__init__(self, plant_name, plant_cost)

        self.is_annual = is_annual
        self.color_of_flowers = color_of_flowers

    def print_info(self):
        Plant.print_info(self)
        print(f"Annual: {self.is_annual}")
        print(f"Color of flowers: {self.color_of_flowers}")


def print_list(my_garden):

    for i in range(len(my_garden)):
        print(f"Plant {i + 1} Information:")
        my_garden[i].print_info()
        print()


my_garden = []

user_string = input()

while user_string != "-1":

    tokens = user_string.split()

    type = tokens[0]
    name = tokens[1]
    cost = tokens[2]

    if type == "plant":
        my_garden.append(Plant(name, cost))

    elif type == "flower":
        is_annual = tokens[3]
        color = tokens[4]
        my_garden.append(Flower(name, cost, is_annual, color))

    user_string = input()

print_list(my_garden)
\end{lstlisting}

\section{Key Takeaways}

\begin{itemize}
    \item A derived class can extend a base class.
    \item A derived class can override a method from the base class.
    \item Lists can store objects of related types.
    \item Polymorphism allows the same method call to behave differently depending on the object.
    \item Input parsing often begins by splitting strings into tokens.
\end{itemize}

\end{document}
